\section{Énoncés des Exercices}

\subsection{Dérivation}

% --- Exercice 1 (Source: EXERCICE 4 analyse) ---
\begin{exercice}[Exercice 4 Analyse]
\begin{enumerate}
    \item Énoncer le théorème des accroissements finis.
    \item Soit $f:[a,b]\longrightarrow\R$ et soit $x_{0}\in]a,b[$. 
        On suppose que $f$ est continue sur $[a, b]$ et que $f$ est dérivable sur $]a, x_{0}[$ et sur $]x_{0}, b[$.
        Démontrer que, si $f'$ admet une limite finie en $x_{0}$, alors $f$ est dérivable en $x_{0}$ et $f'(x_{0})=\lim_{x\rightarrow x_{0}}f'(x)$.
    \item Prouver que l'implication : ($f$ est dérivable en $x_{0}) \Rightarrow$ ($f'$ admet une limite finie en $x_{0}$) est fausse.
        \textit{Indication : on pourra considérer la fonction $g$ définie par : $g(x)=x^{2}\sin\frac{1}{x}$ si $x\ne0$ et $g(0)=0$.}
\end{enumerate}
\end{exercice}

% --- Exercice 2 (Source: Exercice 1) ---
\begin{exercice}
    Étudier (prolongement, variations, courbes) sur $\R^{+}_*$ la fonction $x\mapsto x^{x}$.
\end{exercice}

% --- Exercice 3 (Source: Exercice 2) ---
\begin{exercice}
Soit $f$ définie au voisinage du réel $a$ prenant ses valeurs dans $E$, un espace vectoriel normé de dimension finie, et continue en $a$.
\begin{enumerate}
    \item Comparer la dérivabilité de $f$ en $a$ et l'existence de $\lim_{t\rightarrow0}\frac{f(a+t)-f(a-t)}{2t}$.
    \item Même question avec l'existence de $\lim_{t\rightarrow0}\frac{f(a+2t)-f(a+t)}{t}=\lambda$.
    
    \textit{Indication : on pourra commencer par montrer que pour $t>0$ et $n\in\mathbb{N}^{*}$ on a :}
    $$ \frac{f(a+t)-f(a)}{t} - \lambda = \sum_{k=0}^{n}\frac{1}{2^{k+1}}\left(\frac{f\left(a+\frac{t}{2^{k}}\right)-f\left(a+\frac{t}{2^{k+1}}\right)}{\frac{t}{2^{k+1}}} - \lambda\right) + \frac{f\left(a+\frac{t}{2^{n+1}}\right)-f(a)}{t} - \frac{1}{2^{n+1}}\lambda $$
\end{enumerate}
\end{exercice}

% --- Exercice 4 (Source: Exercice 3) ---
\begin{exercice}
    Étudier la dérivabilité sur $[0, 1]$ de $x\mapsto \arcsin(1-x^{3})$.
\end{exercice}

% --- Exercice 5 (Source: Exercice 4) ---
\begin{exercice}
    Montrer que $t\mapsto e^{-\frac{1}{|t|}}$ est de classe $\mathcal{C}^{\infty}$ sur $\R$.
\end{exercice}

% --- Exercice 6 (Source: Exercice 5) ---
\begin{exercice}
    Soit $A\in \mathcal{M}_{n}(\R)$ et $\varphi\in \mathcal{F}(\R,\R)$ définie par $\forall t\in\R$, $\varphi(t)=\det(tA+I_{n})$.
    Montrer que $\varphi$ est dérivable en 0 et que $\varphi'(0)=\text{Tr}(A)$.
\end{exercice}

% --- Exercice 7 (Source: Exercice 6) ---
\begin{exercice}
    Soit $f$ dérivable sur $[a, b]$ à valeurs dans $\R$.
    \begin{enumerate}
        \item On suppose que $f'(a)>0$ et $f'(b)<0$.
            Montrer qu'il existe $c\in]a,b[$ vérifiant $f'(c)=0$.

        \textit{Indication : on utilisera le fait que $f$ est continue sur le segment et on montrera qu'il n'est pas possible que ses deux extremums ne soient atteints qu'en $a$ et $b$.}
    \item En déduire le théorème de Darboux : "Toute dérivée d'une fonction dérivable de $[a, b]$ dans $\R$ vérifie le théorème des valeurs intermédiaires".
\end{enumerate}
\end{exercice}

% --- Exercice 8 (Source: Exercice 7) ---
\begin{exercice}
    Soit $f\in \mathcal{C}^{1}([a,b],\R)$ telle que $f'(a)=f'(b)$.
    Montrer qu'il existe $c\in]a,b[$ tel que
    $$ f'(c)=\frac{f(c)-f(a)}{c-a} $$
    \textit{Indication : on définira $g:x\mapsto\frac{f(x)-f(a)}{x-a}$ prolongée par continuité sur $[a, b]$ et on montrera qu'elle atteint un de ses extremums en un point $c\in]a,b[$.}
\end{exercice}

% --- Exercice 9 (Source: EXERCICE 3 analyse) ---
\begin{exercice}[Exercice 3 Analyse]
\begin{enumerate}
    \item On pose $g(x)=e^{2x}$ et $h(x)=\frac{1}{1+x}$.
        Calculer, pour tout entier naturel $k$, la dérivée d'ordre $k$ des fonctions $g$ et $h$ sur leurs ensembles de définitions respectifs.
    \item On pose $f(x)=\frac{e^{2x}}{1+x}$.
        En utilisant la formule de Leibniz, concernant la dérivée $n$-ième d'un produit de fonctions, 
        déterminer, pour tout entier naturel $n$ et pour $x\in\R\setminus\{-1\}$, la valeur de $f^{(n)}(x)$.
    \item Démontrer, dans le cas général, la formule de Leibniz, utilisée dans la question précédente.
\end{enumerate}
\end{exercice}

% --- Exercice 10 (Source: Exercice 8) ---
\begin{exercice}
    Calculer la dérivée $n$-ième de $t \mapsto t^{n-1} \ln t$.
    Procéder par récurrence et formule de Leibniz.
\end{exercice}

% --- Exercice 11 (Source: Exercice 9) ---
\begin{exercice}
    Calculer de deux manières différentes la dérivée $n$-ième de $x\mapsto x^{2n}=x^{n}.x^{n}$ 
    pour calculer la valeur de $\displaystyle \sum_{k=0}^{n}\binom{n}{k}^{2}$.
\end{exercice}

\subsection{Intégration sur un segment}

% --- Exercice 12 (Source: EXERCICE 56 analyse) ---
\begin{exercice}[Exercice 56 Analyse]
On considère la fonction $H$ définie sur $]1, +\infty[$ par $\displaystyle H(x)=\int_{x}^{x^{2}}\frac{\dt}{\ln t}$.
\begin{enumerate}
    \item Montrer que $H$ est $\mathcal{C}^{1}$ sur $]1,+\infty[$ et calculer sa dérivée.
    \item Montrer que la fonction $u$ définie par $u(x)=\frac{1}{\ln x}-\frac{1}{x-1}$ admet une limite finie en $x=1$.
    \item En utilisant la fonction $u$ de la question 2, calculer la limite en $1^+$ de la fonction $H$.
\end{enumerate}
\end{exercice}

% --- Exercice 13 (Source: Exercice 10) ---
\begin{exercice}
    Soit $f\in \mathcal{C}^{2}(\R,\R)$ telle que $f+f''>0$. Démontrer que
    $$ \forall x\in\R,~f(x)+f(x+\pi)\ge0 $$
    \textit{Indication : considérer $\displaystyle \int_{0}^{\pi}\sin(t)(f(t)+f''(t))\dt$.}
\end{exercice}

% --- Exercice 14 (Source: Exercice 11) ---
\begin{exercice}
    Déterminer $\displaystyle \lim_{x\rightarrow0^{+}}\int_{x}^{2x}\frac{e^{t}}{\arcsin t}\dt$.
\end{exercice}

% --- Exercice 15 (Source: Exercice 12) ---
\begin{exercice}
    Étudier la fonction définie par $\displaystyle f(x)=\int_{x}^{x^{2}}\frac{\dt}{\ln t}$.
\end{exercice}

% --- Exercice 16 (Source: Exercice 13) ---
\begin{exercice}
    Après avoir montré son existence, en utilisant un changement de variable affine échangeant les bornes, 
    calculer $\displaystyle \int_{0}^{\pi}\frac{x}{1+\sin x}\dx$.
\end{exercice}

% --- Exercice 17 (Source: Exercice 14) ---
\begin{exercice}
    Soit $f$ de classe $\mathcal{C}^{1}$ sur $\R^{+}$, strictement croissante et réalisant une bijection de $\R^{+}$ sur $\R^{+}$.
    Montrer que pour tout $x\geqslant0$,
    $$ xf(x)=\int_{0}^{x}f(t)\dt+\int_{0}^{f(x)}f^{-1}(t)\dt $$
    \textit{Indication : dériver. Interprétation géométrique ?}
\end{exercice}

% --- Exercice 18 (Source: Exercice 15) ---
\begin{exercice}
    Soit $f\in \mathcal{C}^{2}([a,b],E)$ telle que $f(a)=f(b)=0$. Montrer que
    $$ \left\| \int_{a}^{b}f(t)\dt \right\| \leqslant \frac{(b-a)^{3}}{12} \sup_{[a,b]} \|f''(t)\| $$
    Poser $P(t)=\frac{1}{2}(t-a)(b-t)$ et intégrer par parties $\ds \int_{a}^{b}P(t)f''(t)\dt$.
\end{exercice}

% --- Exercice 19 (Source: Exercice 16) ---
\begin{exercice}
    Soit $A\in \mathcal{M}_{n}(\C)$. Calculer de deux manières différentes, pour $r>0$ assez grand,
    $$ I(r)=\int_{0}^{2\pi}re^{i\theta} \det(re^{i\theta}I_{n}-A) (re^{i\theta}I_{n}-A)^{-1} d\theta $$
    Et en déduire le théorème de Cayley-Hamilton.
    \begin{enumerate}
        \item Commencer par montrer que pour $r>0$ assez grand $\left(re^{i\theta}I_{n}-A\right)$ e
            st bien inversible pour tout $\theta$, ce qui assure bien l'existence de $I(r)$.
            On supposera dans toute la suite que $r>0$ est assez grand pour que $I(r)$ existe.
        \item 
        \begin{enumerate}
            \item Utiliser la relation entre inverse et comatrice pour montrer que 
                $\ds I(r)=\int_{0}^{2\pi}re^{i\theta} \text{Com}\left(re^{i\theta}I_n-A\right)^{T} d\theta$.
            \item En déduire que pour tous $i,j\in\{1,\dots,n\}$ il existe un polynôme $Q_{ij}$ tel que 
                $\ds (I(r))_{i,j}=\int_{0}^{2\pi}re^{i\theta}Q_{i,j}\left(re^{i\theta}\right) d\theta$.
            \item En déduire que $I(r)=0$.
        \end{enumerate}
        \item 
        \begin{enumerate}
            \item Montrer que pour tout matrice de norme assez petite, $\ds (I_{n}-B)^{-1}=\sum_{k=0}^{+ \infty}B^{k}$.
            \item Utiliser ce résultat et une interversion somme/intégrale (que l'on admet pour le moment) pour montrer que :
                $$ I(r)=\sum_{k=0}^{+ \infty}r^{-k}A^{k}\int_{0}^{2\pi}\det(re^{i\theta}I_{n}-A)e^{-ik\theta}d\theta $$
                (note: correction dans l'exposant de $r$ selon le contexte classique).
            \item En notant le polynôme caractéristique de $A$, $\ds\chi_{A}=\det(XI_{n}-A)=\sum_{p=0}^{n}a_{p}X^{p}$, 
                conclure en remarquant que $\ds \int_{0}^{2\pi}e^{i(p-k)\theta}d\theta=0$ sauf si $k=p$.
        \end{enumerate}
    \end{enumerate}
\end{exercice}

% --- Exercice 20 (Source: Exercice 17) ---
\begin{exercice}
    Déterminer $\ds\lim_{n\rightarrow+\infty}\sum_{k=1}^{n}\sin\left(\frac{k}{n}\right)\sin\left(\frac{k}{n^{2}}\right)$.
    On se ramènera par un encadrement de $\sin x$ à la recherche de $\ds \lim_{n\rightarrow+\infty}\sum_{k=1}^{n}\sin\left(\frac{k}{n}\right)\frac{k}{n^{2}}$.
\end{exercice}

% --- Exercice 21 (Source: Exercice 18) ---
\begin{exercice}
    Trouver un équivalent de $\ds \sum_{k=1}^{n}k \cos(\frac{k\pi}{n})$.
\end{exercice}

% --- Exercice 22 (Source: Exercice 19) ---
\begin{exercice}
    Soit $x$ réel différent de -1 et 1.
    Après avoir vérifié son existence, utiliser les sommes de Riemann pour calculer $\ds \int_{0}^{2\pi}\ln(x^{2}-2x \cos t + 1)\dt$.
\end{exercice}

\newpage
\section{Corrections}

% --- Correction 2 (Source: Exercice 1) ---
\begin{correction}{2}
    Soit $f : x \mapsto x^x$ définie sur $\R^+_*$.
    On écrit $f(x) = e^{x \ln x}$.

    $f$ est $\mathcal{C}^{\infty}$ par théorèmes d'opérations.
    $\forall x \in \R^+_*, f'(x) = (\ln x + 1) e^{x \ln x}$.

    \begin{itemize}
    \item 
        \textbf{Variations :}
        $f'(x) = 0 \iff \ln x = -1 \iff x = e^{-1} = 1/e$.

        Sur $]0, 1/e[$, $f'(x) < 0$, $f$ décroissante.
        Sur $]1/e, +\infty[$, $f'(x) > 0$, $f$ croissante.

        Minimum en $1/e$ valant $(1/e)^{1/e} = e^{-1/e}$.
    \item
        \textbf{Limites :}
        en $0^+$ : $x \ln x \to 0$, donc $f(x) \to e^0 = 1$.
        On la prolonge par continuité en posant $f(0)=1$.
        
        Dérivabilité en 0 du prolongement :

        \[\frac{f(x)-f(0)}{x} = \frac{e^{x \ln x}-1}{x} \sim_{0} \frac{x \ln x}{x} = \ln x \to -\infty \]
        Donc le prolongement a une tangente verticale en $0$ (et donc n'est pas dérivable en $0$).

        De plus, $f(x) \limit{x}{+\infty} +\infty$.
    \item
        \textbf{Courbe :} branche parabolique de direction $(Oy)$ car $f(x)/x = x^{x-1} \to + \infty$ (si $x>2$).
    \end{itemize}
\end{correction}

% --- Correction 3 (Source: Exercice 2) ---
\begin{correction}{3}
\begin{enumerate}
    \item $\Rightarrow$ On considère le taux $\Delta(t) = \frac{f(a+t)-f(a-t)}{2t}$.

        Si $f$ est dérivable en $a$, $\Delta(t) = \frac{1}{2} (\frac{f(a+t)-f(a)}{t} + \frac{f(a-t)-f(a)}{-t}) \to \frac{1}{2}(f'(a)+f'(a)) = f'(a)$.

        La limite existe et vaut $f'(a)$.

        $\Leftarrow$ Réciproque fausse : Contre-exemple $f(x)=|x|$ en $a=0$.
        $\frac{|t|-|-t|}{2t} = 0 \to 0$, pourtant $f$ non dérivable en 0.

    \item
    $\Rightarrow$ Si $f$ est dérivable par sensiblement le même calcul qu'au dessus le quotient a bien une limite.

    $\Leftarrow$ On démontre la relation via récurrence puis on utilise la définition de limite et la continuité de $f$.
\end{enumerate}
\end{correction}

% --- Correction 4 (Source: Exercice 3) ---
\begin{correction}{4}

    $f(x) = \arcsin(1-x^3)$. Définie si $-1 \le 1-x^3 \le 1 \iff 0 \le x^3 \le 2 \iff 0 \le x \le \sqrt[3]{2}$.

    Sur $[0, 1]$, $f$ est définie et continue.
    Dérivabilité sur $]0, 1]$ : composition de fonctions dérivables.

    $$f'(x) = \frac{-3x^2}{\sqrt{1-(1-x^3)^2}} = \frac{-3x^2}{\sqrt{1-(1-2x^3+x^6)}} = \frac{-3x^2}{\sqrt{2x^3-x^6}} = \frac{-3x^2}{x^{3/2}\sqrt{2-x^3}} = \frac{-3\sqrt{x}}{\sqrt{2-x^3}}$$.

    En $0^+$ : $f'(x) \to 0$. Par le théorème de la limite de la dérivée (puisque $f$ est continue en 0), $f$ est dérivable en 0 et $f'(0)=0$.
\end{correction}

% --- Correction 6 (Source: Exercice 5) ---
\begin{correction}{6}

    Soit $\varphi(t) = \det(tA+I_n)$.
    On sait que $\det(I+H) = 1 + \text{Tr}(H) + o(\|H\|)$.

    $\varphi(t) = \det(I+tA) = 1 + \text{Tr}(tA) + o(t) = 1 + t\text{Tr}(A) + o(t)$.

    Donc $\frac{\varphi(t)-\varphi(0)}{t} = \frac{1+t\text{Tr}(A)+o(t)-1}{t} \to \text{Tr}(A)$.
    Donc $\varphi'(0) = \text{Tr}(A)$.

    Autre méthode : $\varphi(t) = \prod (1+t\lambda_i)$. $\varphi'(t) = \sum \lambda_i \prod_{j \ne i} (1+t\lambda_j)$.
    En $0$, $\varphi'(0) = \sum \lambda_i = \text{Tr}(A)$.
\end{correction}

% --- Correction 7 (Source: Exercice 6) ---
\begin{correction}{7}
\begin{enumerate}
    \item À écrire

    \item Théorème de Darboux : Si $f$ est dérivable, $f'$ vérifie le TVI.
    
        Soit $y$ entre $f'(a)$ et $f'(b)$. On pose $g(x) = yx - f(x)$.

        $g'(a) = y - f'(a)$ et $g'(b)= y - f'(b)$ sont de signes contraires.

        D'après 1., il existe $c$ tel que $g'(c)=0$, soit $f'(c)=y$.
\end{enumerate}
\end{correction}

% --- Correction 8 (Source: Exercice 7) ---
\begin{correction}{8}

    Posons $g(x) = \frac{f(x)-f(a)}{x-a}$ pour $x \ne a$ et $g(a)=f'(a)$.

    $g$ est continue sur $[a,b]$.
    $g(b) = \frac{f(b)-f(a)}{b-a}$.

    Si $g$ atteint un extremum en $c \in ]a,b[$, alors $g'(c)=0$.

    Calculons $g'(x)$ pour $x \ne a$ : $g'(x) = \frac{f'(x)(x-a) - (f(x)-f(a))}{(x-a)^2} = \frac{f'(x)-g(x)}{x-a}$.

    Si $g'(c)=0$, alors $f'(c) = g(c) = \frac{f(c)-f(a)}{c-a}$.

    Reste à montrer que $g$ a un extremum sur $]a,b[$.

    Supposons que $g$ a son maximum en $a$ et son minimum en $b$ (quitte à échanger).
    
    $$g(a) \leqslant g(x) \leqslant g(b)$$

    Par le lemme des trois pentes :
    $$\frac{f(b) - f(a)}{b-a} \leqslant \frac{f(b) - f(x)}{b-x}$$

    On fait tendre $x$ vers $b$ :
    $$\frac{f(b) - f(a)}{b-a} \leqslant f'(b) = f'(a)$$

    Donc $\forall x \in ]a,b[, \frac{f(x) -f(a)}{x-a} = f'(a)$.
    Donc $f$ affine.
\end{correction}

\begin{correction}{10}
Notons $f_n: t \to t^{n-1} \ln(t)$ et $g_n: t \to t^{n-1}$.    
Les deux fonctions sont bien $\mathscr C^n$ donc $f$ l'est aussi.
On a :
\[\forall K \in \inter 0 {n-1}, \left(g_n\right)^{(K)}(x) = \frac{(n-1)!}{(n-1-K)!}x^{n-1-K}\]
Et :
\[\forall K \in \inter 1 {n}, \left(\ln\right)^{(K)}(x) = \frac{(-1)^K(K-1)!}{x^K}\]
On calcule avec la formule de Leibniz :
\begin{align*}
    \left(f_n\right)^{(n)} &= \sum_{k=0}^n \binom{n}{k} g_n^{(n-k)}(x) \ln^{k}(x) \\
    &= \sum_{k=1}^n \binom{n}{k} \frac{(n-1)!}{(k-1)!} x^{k-1} \frac{(-1)^{k+1}(k-1)!}{x^k} \\
    &= \frac{(n-1)!}{x} \sum_{k=1}^n \binom{n}{k} (-1)^{k+1} \\
    &= \frac{(n-1)!}{x} \sum_{k=0}^n \left(\binom{n}{k} (-1)^{k+1} -1\right)\\
    &= \frac{(n-1)!}{x}
\end{align*}
\end{correction}

% --- Correction 11 (Source: Exercice 9) ---
\begin{correction}{11}
    On note $f_n: x \to x^n$.

    Dérivée $n$-ième : $(f_{2n})^{(n)}(x) = \frac{(2n)!}{(2n-n)!} x^{2n-n} = \frac{(2n)!}{n!} x^n$.
    En $x=1$, valeur $\frac{(2n)!}{n!}$. \\
    
    Leibniz sur $x^n \cdot x^n$.
    $\ds(f_{2n})^{(n)}(x) = \sum_{k=0}^n \binom{n}{k} (f_n)^{(k)}(x) (f_n)^{(n-k)}(x)$.

    Or $(f_n)^{(k)}(x) = \frac{n!}{(n-k)!} x^{n-k}$.
    En $x=1$, $(f_n)^{(k)}(1) = \frac{n!}{(n-k)!}$.

    Donc $\ds (f_{2n})^{(n)}(1) = \sum_{k=0}^n \binom{n}{k} \frac{n!}{(n-k)!} \frac{n!}{k!} = n! \sum_{k=0}^n \binom{n}{k} \binom{n}{n-k} = n! \sum_{k=0}^n \binom{n}{k}^2$. \\

    Égalité des résultats : $\ds \frac{(2n)!}{n!} = n! \sum_{k=0}^n \binom{n}{k}^2$.
    D'où $\ds \sum_{k=0}^n \binom{n}{k}^2 = \frac{(2n)!}{(n!)^2} = \binom{2n}{n}$.
\end{correction}

% --- Correction 13 (Source: Exercice 10) ---
\begin{correction}{13}
    Posons $\ds I = \int_0^\pi (f(t)+f''(t)) \sin t \dif t$.
    Par IPP (deux fois) sur $\int f'' \sin$ :
    \begin{align*}
        \int_0^\pi f''(t) \sin t \dif t &= [f'(t)\sin t]_0^\pi - \int_0^\pi f'(t) \cos t \dt \\
        &= 0 - \left([f(t)\cos t]_0^\pi - \int_0^\pi f(t)(-\sin t)\dt\right) \\
        &= - \left(-f(\pi) - f(0) + \int_0^\pi f(t) \sin t \dt\right) \\
        &= f(\pi)+f(0) - \int_0^\pi f(t) \sin t \dt
    \end{align*}
    
    Donc $\ds I = \int_0^{\pi} f \sin + f(\pi)+f(0) - \int_0^{\pi} f \sin = f(\pi)+f(0)$.

    Comme $f+f''>0$ et $\sin t > 0$ sur $]0,\pi[$, l'intégrale $I$ est strictement positive.
    Donc $f(\pi)+f(0) > 0$.

    En appliquant ce résultat à $x \mapsto f(x+t)$, on montre que pour tout $t$, $f(t+\pi)+f(t) > 0$.
\end{correction}

% --- Correction 14 (Source: Exercice 11) ---
\begin{correction}{14}

    Soit $\ds f(x) = \int_x^{2x} \frac{e^t}{\arcsin t} \dt$.
    Quand $t \to 0$, $\frac{e^t}{\arcsin t} \sim \frac{1}{t}$.

    On compare avec $\ds g(x) = \int_x^{2x} \frac{1}{t} \dt = [\ln t]_x^{2x} = \ln 2$.
    $\ds f(x) - g(x) = \int_x^{2x} \left(\frac{e^t}{\arcsin t} - \frac{1}{t}\right) \dt$

    La fonction $h(t) = \frac{e^t}{\arcsin t} - \frac{1}{t}$ se prolonge par continuité en 0 (limite 1).
    
    Donc $\ds \int_x^{2x} h(t) \dif t \approx x \cdot h(0) \to 0$.

    Ainsi, $\lim_{x\to 0^+} f(x) = \ln 2$.
\end{correction}

% --- Correction 16 (Source: Exercice 13) ---
\begin{correction}{16}

    $\ds I = \int_0^\pi \frac{x}{1+\sin x} \dx$ existe bien car ce qui est intégré est continu sur $[0,\pi]$.

    Changement de variable $u = \pi - x$. $\dx = -\dif u$.
    Bornes : $\pi \to 0$.
    
    $$I = \int_\pi^0 \frac{\pi-u}{1+\sin(\pi-u)} (-\dif u) = \int_0^\pi \frac{\pi-u}{1+\sin u} \dif u$$

    $$2I = \int_0^\pi \frac{x + \pi - x}{1+\sin x} \dx = \pi \int_0^\pi \frac{1}{1+\sin x} \dx$$

    Calcul de $J = \int_0^\pi \frac{1}{1+\sin x} \dx$.
    On peut utiliser $1+\sin x = 1 + \cos(\frac{\pi}{2}-x) = 2 \cos^2(\frac{\pi}{4}-\frac{x}{2})$.
    \begin{align*}
        J &= \frac{1}{2} \int_0^\pi \frac{\dif x}{\cos^2(\frac{\pi}{4}-\frac{x}{2})} &\text{primitive} \tan \\
          &= [\tan(\frac{x}{2}-\frac{\pi}{4})]_0^\pi \\
          &= \tan(\frac{\pi}{4}) - \tan(-\frac{\pi}{4}) \\
          &= 1 - (-1) \\
          &= 2
    \end{align*}
    Donc $2I = 2\pi \implies I = \pi$.
\end{correction}

% --- Correction 20 (Source: Exercice 17) ---
\begin{correction}{20}

    On cherche $\lim \sum \sin(k/n) \sin(k/n^2)$.

    Pour $n$ grand, $k/n^2$ est petit. $\sin(k/n^2) \sim k/n^2$.

    On compare la somme à $S_n = \sum_{k=1}^n \sin(k/n) \frac{k}{n^2} = \frac{1}{n} \sum_{k=1}^n \frac{k}{n} \sin(\frac{k}{n})$.

    C'est une somme de Riemann pour la fonction $x \mapsto x \sin x$ sur $[0,1]$.
    Elle converge vers $\int_0^1 x \sin x \dx$.

    IPP : $[-x\cos x]_0^1 + \int_0^1 \cos x \dx = -\cos 1 + [\sin x]_0^1 = \sin 1 - \cos 1$.

    Majoration de l'erreur : $|\sin u - u| \le \frac{|u|^3}{6}$.

    L'erreur totale tend vers 0. La limite est $\sin 1 - \cos 1$.
\end{correction}

% --- Correction 21 (Source: Exercice 18) ---
\begin{correction}{21}
    $S_n = \sum_{k=1}^n k \cos(\frac{k\pi}{n})$.

    On écrit $S_n = n^2 \frac{1}{n} \sum_{k=1}^n \frac{k}{n} \cos(\frac{k\pi}{n})$.

    La somme est une somme de Riemann pour $x \mapsto x \cos(\pi x)$ sur $[0,1]$.

    L'intégrale est $\int_0^1 t \cos(\pi t) \dt = \dots = -2/\pi^2$.

    Donc $S_n \sim n^2 \times (-2/\pi^2)$.
\end{correction}

% --- Correction 22 (Source: Exercice 19) ---
\begin{correction}{22}
    $I = \int_0^{2\pi} \ln(x^2 - 2x \cos t + 1) \dt$.

    Sommes de Riemann : 
    \begin{math}
        R_n = \frac{2\pi}{n} \sum_{k=0}^{n-1} \ln(x^2 - 2x \cos(2k\pi/n) + 1) \\
        x^2 - 2x \cos \theta + 1 = |x - e^{i\theta}|^2 \\
        \sum \ln |x-e^{i\theta_k}|^2 = 2 \ln \prod |x-e^{2ik\pi/n}| = 2 \ln |x^n-1| \\
        R_n = \frac{4\pi}{n} \ln |x^n-1|
    \end{math}
    \begin{itemize}
    \item 
        Si $|x| > 1$, $|x^n-1| \sim |x|^n$. $R_n \sim \frac{4\pi}{n} n \ln |x| = 4\pi \ln |x|$.
    \item
        Si $|x| < 1$, $|x^n-1| \to 1$. $R_n \to 0$.
    \end{itemize}
    Donc $I = 4\pi \ln \max(1, |x|)$.
\end{correction}
